\documentclass[11pt]{article}
\usepackage{amsmath, amssymb, amsthm}
\usepackage{geometry}
\usepackage{hyperref}
\usepackage{graphicx}
\usepackage{booktabs}

\geometry{margin=1in}

\title{EE/CS 148B HW4: Diffusion Models}
\author{Dhruv Sheth}
\date{}

\begin{document}
\maketitle

\section*{Part 1: DDPM}

\subsection*{Problem 1.2 (Variational Bound)}

We minimize $-\log p_\theta(x_0)$ to maximize likelihood. Computing $p_\theta(x_0)$ requires marginalizing over all latents $x_1, \ldots, x_T$, which is intractable. Instead we minimize the ELBO, which upper-bounds $-\log p_\theta(x_0)$ by Jensen's inequality applied to the concave log function:
\[
-\log p_\theta(x_0) \leq \mathbb{E}_q\left[-\log \frac{p_\theta(x_{0:T})}{q(x_{1:T}|x_0)}\right].
\]
Specifically, $-\log p_\theta(x_0) = -\log \mathbb{E}_q\left[\frac{p_\theta(x_{0:T})}{q(x_{1:T}|x_0)}\right] \leq \mathbb{E}_q\left[-\log \frac{p_\theta(x_{0:T})}{q(x_{1:T}|x_0)}\right]$, where the inequality is Jensen's. The chain of inequalities in Eq.~(3) follows by expanding the joint and Markov structure.

\subsection*{Problem 1.3 (Tractable Loss)}

\textbf{Eq.~(18):} The posterior $q(x_{t-1}|x_t, x_0)$ has a closed-form Gaussian because $q(x_t|x_{t-1})$ and $q(x_{t-1}|x_0)$ are both Gaussian, and the product of Gaussian densities is proportional to a Gaussian.

\textbf{Eq.~(20):} When both distributions have the same variance, the KL between two Gaussians reduces to $\frac{1}{2\sigma^2}\|\mu_1 - \mu_2\|^2$ by the standard closed-form KL formula.

\textbf{Eq.~(21):} Reparametrizing $x_t = \sqrt{\bar\alpha_t}\, x_0 + \sqrt{1-\bar\alpha_t}\,\varepsilon$ with $\varepsilon \sim \mathcal{N}(0, I)$ lets us compute the loss without simulating the full Markov chain.

\textbf{Eq.~(22):} Substituting the optimal $\mu_\theta$ parameterization from Eq.~(11) into Eq.~(10) and simplifying yields Eq.~(22).

\subsection*{Problem 1.4}

\textbf{Claim:} $q(x_t | x_0) = \mathcal{N}\!\left(\sqrt{\bar\alpha_t}\, x_0,\; (1 - \bar\alpha_t) I\right)$, where $\bar\alpha_t = \prod_{i=1}^t (1-\beta_i)$.

\begin{proof}
By induction on $t$.

\textbf{Base case ($t=1$):} By definition $q(x_1|x_0) = \mathcal{N}(\sqrt{1-\beta_1}\,x_0, \beta_1 I)$. Since $\alpha_1 = 1-\beta_1$ and $\bar\alpha_1 = \alpha_1$, this equals $\mathcal{N}(\sqrt{\bar\alpha_1}\,x_0, (1-\bar\alpha_1)I)$.

\textbf{Inductive step:} Assume $q(x_{t-1}|x_0) = \mathcal{N}(\sqrt{\bar\alpha_{t-1}}\,x_0, (1-\bar\alpha_{t-1})I)$. The transition is $x_t = \sqrt{1-\beta_t}\,x_{t-1} + \sqrt{\beta_t}\,\varepsilon_t$ with $\varepsilon_t \sim \mathcal{N}(0,I)$ independent of $x_{t-1}$. Then:
\begin{align*}
\mathbb{E}[x_t | x_0] &= \sqrt{1-\beta_t}\cdot\sqrt{\bar\alpha_{t-1}}\,x_0 = \sqrt{\alpha_t \bar\alpha_{t-1}}\,x_0 = \sqrt{\bar\alpha_t}\,x_0, \\
\mathrm{Var}[x_t | x_0] &= (1-\beta_t)(1-\bar\alpha_{t-1}) + \beta_t = 1 - \bar\alpha_{t-1} + \bar\alpha_{t-1}\beta_t - \beta_t + \beta_t \\
&= 1 - (1-\beta_t)\bar\alpha_{t-1} = 1 - \alpha_t\bar\alpha_{t-1} = 1 - \bar\alpha_t.
\end{align*}
Since $x_t$ is an affine transformation of jointly Gaussian variables, $q(x_t|x_0) = \mathcal{N}(\sqrt{\bar\alpha_t}\,x_0, (1-\bar\alpha_t)I)$.
\end{proof}

\subsection*{Problem 1.5}

\textbf{Claim:} $q(x_{t-1}|x_t, x_0)$ is Gaussian with mean and variance given by Eqs.~(6) and (7).

\begin{proof}
By Bayes' rule and the Markov property:
\[
q(x_{t-1}|x_t, x_0) \propto q(x_t|x_{t-1})\, q(x_{t-1}|x_0).
\]
Both factors are Gaussian. Let $q(x_t|x_{t-1}) = \mathcal{N}(\sqrt{\alpha_t}\,x_{t-1}, \beta_t I)$ and $q(x_{t-1}|x_0) = \mathcal{N}(\sqrt{\bar\alpha_{t-1}}\,x_0, (1-\bar\alpha_{t-1})I)$. The product of two Gaussian densities in $x_{t-1}$ is proportional to a Gaussian with precision and mean:
\[
\frac{1}{\tilde\sigma_t^2} = \frac{\alpha_t}{\beta_t} + \frac{1}{1-\bar\alpha_{t-1}}, \qquad
\frac{\tilde\mu_t}{\tilde\sigma_t^2} = \frac{\sqrt{\alpha_t}}{\beta_t}\,x_t + \frac{\sqrt{\bar\alpha_{t-1}}}{1-\bar\alpha_{t-1}}\,x_0.
\]
Computing $\tilde\sigma_t^2$:
\[
\frac{1}{\tilde\sigma_t^2} = \frac{\alpha_t(1-\bar\alpha_{t-1}) + \beta_t}{\beta_t(1-\bar\alpha_{t-1})} = \frac{1-\bar\alpha_t}{\beta_t(1-\bar\alpha_{t-1})},
\]
so
\[
\tilde\sigma_t^2 = \frac{\beta_t(1-\bar\alpha_{t-1})}{1-\bar\alpha_t}.
\]
Multiplying through for $\tilde\mu_t$:
\[
\tilde\mu_t = \frac{\sqrt{\alpha_t}(1-\bar\alpha_{t-1})}{1-\bar\alpha_t}\,x_t + \frac{\sqrt{\bar\alpha_{t-1}}\,\beta_t}{1-\bar\alpha_t}\,x_0,
\]
which are Eqs.~(6) and (7).
\end{proof}

\subsection*{Problem 1.6}

The KL between two Gaussians $\mathcal{N}(\mu_1, \sigma^2 I)$ and $\mathcal{N}(\mu_2, \sigma^2 I)$ with equal diagonal covariance $\sigma^2 I$ is:
\[
D_\mathrm{KL}(\mathcal{N}(\mu_1, \sigma^2 I) \| \mathcal{N}(\mu_2, \sigma^2 I)) = \frac{1}{2\sigma^2}\|\mu_1 - \mu_2\|^2.
\]
In the DDPM loss, $L_{t-1} = D_\mathrm{KL}(q(x_{t-1}|x_t, x_0) \| p_\theta(x_{t-1}|x_t))$. Taking $q(x_{t-1}|x_t,x_0) = \mathcal{N}(\tilde\mu_t, \tilde\sigma_t^2 I)$ and $p_\theta(x_{t-1}|x_t) = \mathcal{N}(\mu_\theta(x_t,t), \tilde\sigma_t^2 I)$ (same variance), we get:
\[
L_{t-1} = \frac{1}{2\tilde\sigma_t^2}\mathbb{E}\!\left[\|\tilde\mu_t(x_t, x_0) - \mu_\theta(x_t, t)\|^2\right].
\]
Setting $\sigma_t^2 = \tilde\sigma_t^2 = \beta_t$ (as in the paper's simplified setting) gives:
\[
L_{t-1} - C = \mathbb{E}\left[\frac{1}{2\beta_t}\|\tilde\mu_t - \mu_\theta\|^2\right],
\]
where $C$ absorbs terms that do not depend on $\theta$.

\subsection*{Problem 1.7}

From Eq.~(11), $\mu_\theta(x_t, t) = \frac{1}{\sqrt{\alpha_t}}\left(x_t - \frac{\beta_t}{\sqrt{1-\bar\alpha_t}}\varepsilon_\theta(x_t, t)\right)$ and the true posterior mean is $\tilde\mu_t = \frac{1}{\sqrt{\alpha_t}}\left(x_t - \frac{\beta_t}{\sqrt{1-\bar\alpha_t}}\varepsilon\right)$. Substituting into Eq.~(10):
\begin{align*}
L_{t-1} - C &= \mathbb{E}\left[\frac{1}{2\tilde\sigma_t^2}\|\tilde\mu_t - \mu_\theta\|^2\right] \\
&= \mathbb{E}\left[\frac{1}{2\tilde\sigma_t^2}\cdot\frac{\beta_t^2}{\alpha_t(1-\bar\alpha_t)}\|\varepsilon - \varepsilon_\theta(\sqrt{\bar\alpha_t}\,x_0 + \sqrt{1-\bar\alpha_t}\,\varepsilon,\, t)\|^2\right] \\
&= \mathbb{E}\left[\frac{\beta_t^2}{2\sigma_t^2 \alpha_t(1-\bar\alpha_t)}\|\varepsilon - \varepsilon_\theta(\sqrt{\bar\alpha_t}\,x_0 + \sqrt{1-\bar\alpha_t}\,\varepsilon,\, t)\|^2\right],
\end{align*}
which is Eq.~(12).

\subsection*{Problem 1.8 (Coefficient Plot)}

See submitted plot. The weighting coefficient $\frac{\beta_t^2}{2\sigma_t^2 \alpha_t(1-\bar\alpha_t)}$ is larger for small $t$ (near the end of denoising, fine-grained steps) and smaller for large $t$. The simplified objective drops these coefficients, effectively down-weighting the loss at small $t$. This has been empirically shown to improve sample quality, as it prevents the model from over-optimizing the easy, low-noise denoising steps at the expense of learning the harder, high-noise steps.

\section*{Part 2: Score Matching}

\subsection*{Problem 2.2}

The explicit score matching objective requires computing $\nabla_x \log p(x)$, the score of the true data distribution. For complex distributions $p(x)$, this is intractable since we do not have an analytic form for $p(x)$, only samples from it.

\subsection*{Problem 2.3}

\textbf{Claim:} $J_\mathrm{ESM}(\theta) = J_\mathrm{ISM}(\theta) + C$ where $C$ does not depend on $\theta$.

\begin{proof}
Expand the ESM objective:
\[
J_\mathrm{ESM}(\theta) = \mathbb{E}_p\!\left[\|s_\theta(x) - \nabla_x \log p(x)\|^2\right] = \mathbb{E}_p\!\left[\|s_\theta\|^2\right] - 2\,\mathbb{E}_p\!\left[s_\theta \cdot \nabla_x \log p\right] + \underbrace{\mathbb{E}_p\!\left[\|\nabla_x \log p\|^2\right]}_{=\,C}.
\]
For the cross term, note $\nabla_x \log p(x) = \nabla_x p(x)/p(x)$, so:
\[
\mathbb{E}_p[s_\theta \cdot \nabla_x \log p] = \int s_\theta(x)\, \nabla_x p(x)\, dx.
\]
Integrating by parts (boundary terms vanish as $p(x) \to 0$ at infinity):
\[
\int s_\theta(x)\, \nabla_x p(x)\, dx = -\int (\nabla_x \cdot s_\theta(x))\, p(x)\, dx = -\mathbb{E}_p[\mathrm{tr}(\nabla_x s_\theta(x))].
\]
Therefore:
\[
J_\mathrm{ESM}(\theta) = \mathbb{E}_p\!\left[\mathrm{tr}(\nabla_x s_\theta(x)) + \tfrac{1}{2}\|s_\theta(x)\|^2\right] + C = J_\mathrm{ISM}(\theta) + C.
\]
\end{proof}

\subsection*{Problem 2.4}

The ISM objective requires computing $\mathrm{tr}(\nabla_x s_\theta(x))$, the trace of the Jacobian of $s_\theta$. Naively this requires $O(d)$ backward passes (one per output dimension), making it $O(d)$ times more expensive than a single forward pass. For high-dimensional data (e.g., $d = 784$ for $28\times28$ images), this is prohibitively expensive.

\subsection*{Problem 2.5 (Tweedie's Formula)}

\textbf{Claim:} For $\tilde{x} = x + \sigma \varepsilon$, $\varepsilon \sim \mathcal{N}(0,I)$, we have $\mathbb{E}[x|\tilde{x}] = \tilde{x} + \sigma^2 \nabla_{\tilde{x}} \log p_\sigma(\tilde{x})$.

\begin{proof}
The marginal density is $p_\sigma(\tilde{x}) = \int p(x)\, \phi_\sigma(\tilde{x}-x)\, dx$ where $\phi_\sigma(u) = (2\pi\sigma^2)^{-d/2}\exp(-\|u\|^2/(2\sigma^2))$.

Taking the gradient:
\[
\nabla_{\tilde{x}} \log p_\sigma(\tilde{x}) = \frac{\nabla_{\tilde{x}} p_\sigma(\tilde{x})}{p_\sigma(\tilde{x})} = \frac{\int p(x)\, \nabla_{\tilde{x}} \phi_\sigma(\tilde{x}-x)\, dx}{p_\sigma(\tilde{x})}.
\]
Since $\nabla_{\tilde{x}} \phi_\sigma(\tilde{x}-x) = -\frac{\tilde{x}-x}{\sigma^2}\phi_\sigma(\tilde{x}-x)$, we get:
\[
\nabla_{\tilde{x}} \log p_\sigma(\tilde{x}) = \frac{\int p(x)\cdot\left(-\frac{\tilde{x}-x}{\sigma^2}\right)\phi_\sigma(\tilde{x}-x)\,dx}{p_\sigma(\tilde{x})} = -\frac{\tilde{x} - \mathbb{E}[x|\tilde{x}]}{\sigma^2}.
\]
Rearranging: $\mathbb{E}[x|\tilde{x}] = \tilde{x} + \sigma^2 \nabla_{\tilde{x}} \log p_\sigma(\tilde{x})$.
\end{proof}

\subsection*{Problem 2.6 (DSM)}

The Tweedie-based training objective minimizes:
\[
\mathbb{E}_{x, \tilde{x}}\!\left[\left\|s_\theta(\tilde{x}) - \nabla_{\tilde{x}} \log p(\tilde{x}|x)\right\|^2\right] = \mathbb{E}_{x,\varepsilon}\!\left[\left\|s_\theta(x + \sigma\varepsilon) + \frac{\varepsilon}{\sigma}\right\|^2\right],
\]
since $\nabla_{\tilde{x}} \log p(\tilde{x}|x) = -(\tilde{x}-x)/\sigma^2 = -\varepsilon/\sigma$. Expanding:
\[
= \mathbb{E}\!\left[\|s_\theta(\tilde{x})\|^2 + \frac{2}{\sigma}s_\theta(\tilde{x})\cdot\varepsilon + \frac{\|\varepsilon\|^2}{\sigma^2}\right].
\]
The last term is constant in $\theta$, so this equals the DSM objective up to a constant independent of $\theta$.

\subsection*{Problem 2.7 (DSM vs DDPM)}

The DDPM noise predictor $\varepsilon_\theta$ and the DSM score $s_\theta$ are related by $s_\theta(x_t, t) = -\varepsilon_\theta(x_t, t)/\sigma_t$, where $\sigma_t = \sqrt{1-\bar\alpha_t}$.

The DSM objective is $\mathbb{E}\|s_\theta(x_t) - \nabla_{x_t} \log p_{\sigma_t}(x_t)\|^2$. Substituting $s_\theta = -\varepsilon_\theta/\sigma_t$ and $\nabla_{x_t}\log p_{\sigma_t}(x_t) = -\varepsilon/\sigma_t$:
\[
\mathbb{E}\left\|\frac{-\varepsilon_\theta}{\sigma_t} - \frac{-\varepsilon}{\sigma_t}\right\|^2 = \frac{1}{\sigma_t^2}\mathbb{E}\|\varepsilon_\theta - \varepsilon\|^2,
\]
which matches the DDPM noise prediction objective up to the constant $c(\sigma_t) = 1/\sigma_t^2$.

\subsection*{Problem 2.8 (Langevin Dynamics)}

To sample from $p(x)$ using a learned score $s_\theta$, run Langevin dynamics: initialize $x_0 \sim q(x)$ (e.g., Gaussian), then iterate
\[
x_{k+1} = x_k + \frac{\delta}{2}\, s_\theta(x_k) + \sqrt{\delta}\, z_k, \quad z_k \sim \mathcal{N}(0, I),
\]
for $K$ steps with step size $\delta > 0$. For small $\sigma$, $p_\sigma \approx p(x)$, so using the score learned under noise level $\sigma$ gives samples approximately from $p(x)$. After enough steps, the chain converges to $p_\sigma \approx p(x)$.

\section*{Part 3: Score-Based SDEs}

\subsection*{Problem 3.2}

Two approximations used to recover the DDPM forward process from the VP-SDE discretization:

(1) $e^{-\frac{1}{2}\beta_i\,\Delta t} \approx 1 - \frac{1}{2}\beta_i\,\Delta t$ by first-order Taylor expansion, valid when $\Delta t \ll 1$.

(2) $(1 - \beta_i\,\Delta t)^{1/2} \approx 1 - \frac{1}{2}\beta_i\,\Delta t$ by the same first-order approximation.

\subsection*{Problem 3.3}

By Anderson (1982), the reverse SDE for the VP-SDE is:
\[
dx = \left[-\frac{1}{2}\beta(t)\, x - \beta(t)\,\nabla_x \log p_t(x)\right]dt + \sqrt{\beta(t)}\, d\bar{B}_t,
\]
where $\bar{B}_t$ is a standard Brownian motion running backward in time.

\subsection*{Problem 3.4}

Discretize the reverse VP-SDE with step size $\Delta t = 1/T$ at time $t = (T-i)/T$. Set $\beta(t)\Delta t = \beta_i$. Approximating $e^{-\beta_i/2} \approx \sqrt{\alpha_i}$ and using $\nabla_x \log p_t(x_t) \approx -\varepsilon_\theta(x_t, t)/\sqrt{1-\bar\alpha_t}$, the discretization becomes:
\[
x_{t-1} = \frac{1}{\sqrt{\alpha_t}}\left(x_t - \frac{1-\alpha_t}{\sqrt{1-\bar\alpha_t}}\varepsilon_\theta(x_t,t)\right) + \sqrt{\beta_t}\, z,\quad z \sim \mathcal{N}(0,I),
\]
which is exactly the DDPM ancestral sampling update.

\subsection*{Problem 3.5}

The coefficient difference between the SDE derivation and DDPM arises from the relationship $s_\theta = -\varepsilon_\theta/\sigma_t$. In the reverse SDE, $\nabla_x \log p_t(x)$ appears directly. Substituting $\nabla_x \log p_t \approx -\varepsilon_\theta/\sqrt{1-\bar\alpha_t}$ into the SDE and simplifying using $1-\alpha_t = \beta_t$ recovers the DDPM coefficient $\beta_t/\sqrt{1-\bar\alpha_t}$.

\subsection*{Problem 3.6}

The reverse SDE uses a time-varying diffusion coefficient $g(t) = \sqrt{\beta(t)}$ and a time-conditioned score $s_\theta(x_t, t)$ that encodes the full trajectory structure. Langevin dynamics uses a fixed noise level $\sigma$ and a fixed score, without exploiting the time structure.

In practice, this means DDPM trains a single network $\varepsilon_\theta(x_t, t)$ conditioned on time, whereas standard Langevin would require a separate model for each noise level. The SDE framework gives a principled way to schedule noise levels and leverage the continuous-time structure for sampling.

\subsection*{Problem 3.7}

By Bayes' rule: $\nabla_x \log p_t(x|y) = \nabla_x \log p_t(x) + \nabla_x \log p_t(y|x)$.

For classifier guidance, train a time-conditional classifier $p_\phi(y|x_t, t)$ and substitute $\nabla_x \log p_t(y|x) \approx \nabla_x \log p_\phi(y|x_t, t)$. Then run the reverse SDE with the combined score:
\[
\tilde{s}(x_t, t) = s_\theta(x_t, t) + w\,\nabla_{x_t} \log p_\phi(y|x_t, t),
\]
where $w \geq 1$ is a guidance scale. This steers the samples toward class $y$ while maintaining image quality.

\section*{Part 4: Rectified Flow}

\subsection*{Problem 4.1}

Differentiating the rectified flow objective $L_\mathrm{RF}(\theta) = \mathbb{E}\|v_\theta(X_t, t) - (X_1 - X_0)\|^2$ over the distribution of $(X_0, X_1, t)$ and setting to zero, the minimizer satisfies:
\[
v^*(x, t) = \mathbb{E}[X_1 - X_0 \mid X_t = x].
\]
This follows because the MSE objective is minimized pointwise by the conditional mean.

\subsection*{Problem 4.2}

Since $X_t = (1-t)X_0 + tX_1$, we have $X_1 - X_0 = (X_t - X_0)/(t) + X_0 - X_0 \cdot \ldots$. More directly, $X_t - X_0 = t(X_1 - X_0)$, so $(X_1 - X_0) = (X_t - X_0)/t$. Then:
\[
v^*(x, t) = \frac{x - \mathbb{E}[X_0 | X_t = x]}{t}.
\]
By Tweedie's formula applied to the marginal $p_t(x)$ (with $\sigma_t^2 = t^2(1-\sigma_{\min}^2)$ in the appropriate noise schedule), $\mathbb{E}[X_0|X_t=x] = x + \sigma_t^2\,\nabla_x \log p_t(x)$. Substituting connects $v^*$ to the score function $\nabla_x \log p_t(x)$.

\subsection*{Problem 4.3}

$S=1$ means perfectly straight trajectories (constant velocity field); $S=0$ means maximally curved paths. Reflow straightens trajectories by re-training $v_\theta$ on $(X_0', X_1')$ pairs generated by running the learned ODE, so the new pairs are already coupled by the flow. These pairs have fewer crossings, and re-training on them produces a straighter velocity field. Each round of reflow has diminishing returns, so 1--2 rounds typically suffice. The cost is one additional training run plus generating approximately 50K pairs, comparable to one training run.

\subsection*{Problem 4.4}

\begin{enumerate}
    \item \textbf{Training targets:} DDPM regresses the noise $\varepsilon$ added at step $t$. Rectified flow regresses the velocity $X_1 - X_0$. The two targets differ by a $\sigma_t$ scaling, since $\varepsilon \propto \sigma_t (X_1 - X_0)$ under the corresponding noise schedule.

    \item \textbf{Inference:} DDPM requires approximately 1000 steps because it follows a stochastic SDE trajectory. After reflow, rectified flow can generate samples in a single step since the ODE trajectories are straight and the Euler method is exact for linear flows.

    \item \textbf{Geometry:} DDPM follows a curved SDE path through data space, requiring small step sizes to control discretization error. Rectified flow follows straight lines, which have smaller discretization error and allow fewer function evaluations.

    \item \textbf{Likelihood:} DDPM provides the ELBO as a variational lower bound on $\log p(x_0)$. Rectified flow is a continuous normalizing flow, so exact likelihood is computable via the instantaneous change-of-variables formula (change-of-variables theorem for ODEs), though evaluating the log-determinant of the Jacobian is expensive without special structure.
\end{enumerate}

\section*{Part 5: DDPM Implementation}

Training was run on a RunPod A100-SXM 80GB instance for 50 epochs on FashionMNIST.
Code is in the submitted \texttt{diffusion/vp.py} and \texttt{scripts/train\_vp.py}.

\begin{figure}[h]
\centering
\includegraphics[width=0.72\textwidth]{ddpm_training_loss.pdf}
\caption{VP-SDE training and validation DSM loss over 50 epochs on FashionMNIST.}
\end{figure}

\section*{Part 6: Score-Based and VP-SDE Sampling}

Training was run on a RunPod A100-SXM 80GB instance. The KID scores across sampling steps are reported in Table~\ref{tab:kid}.

\begin{table}[h]
\centering
\begin{tabular}{lrrrrrrr}
\toprule
Method & 1 step & 5 steps & 10 steps & 50 steps & 100 steps & 200 steps & 1000 steps \\
\midrule
VP EM     & 0.5517 & 0.3990 & 0.3334 & 0.1187 & 0.0760 & 0.0600 & 0.0512 \\
RectFlow  & 0.3845 & 0.0263 & 0.0131 & 0.0060 & 0.0067 & 0.0049 & 0.0065 \\
Reflow    & 0.0456 & 0.0157 & 0.0156 & 0.0141 & 0.0137 & 0.0134 & 0.0123 \\
\bottomrule
\end{tabular}
\caption{KID scores for VP Euler--Maruyama, RectFlow, and Reflow at various sampling steps.}
\label{tab:kid}
\end{table}

The results show that RectFlow significantly outperforms VP EM at intermediate step counts (5--50 steps), and Reflow achieves the best 1-step KID by a large margin (0.0456 vs 0.3845 for RectFlow), demonstrating the effectiveness of trajectory straightening for few-step generation.

\section*{Part 7: Rectified Flow}

Code is in \texttt{diffusion/rectflow.py} and \texttt{scripts/train\_rectflow.py}.

\end{document}
